\section{Introduction and research question}\label{sec:introduction}
Text-conditioned image generation is informative only when the conditioning request can be evaluated separately from generic image plausibility. Here, each caption specifies face shape, glasses, hair colour, and face colour. The generator must map this structured request to visible attributes and, more demanding still, generalize to attribute combinations that were deliberately excluded from training.

The final deliverable is a four-model $2\times2$ family: unconditional and conditional endpoints at native32/40k and native64/60k. The scientific questions are deliberately narrower than a universal model ranking. We characterize structured prompt adherence for the conditional endpoints, marginal image quality and diversity for the unconditional endpoints, and a five-endpoint conditional ablation sequence. The CA8-to-H7 comparison is informative but not a pure attention-placement estimate because it combines independent trajectories with a parameter increase; V2-A is a native64 extension, not a resolution-only ablation. This framing prevents observed differences from being attributed to a single factor when resolution, preprocessing, capacity, or training trajectory also changed.

Development evidence is separated from confirmatory evaluation. It can motivate a protocol and diagnose implementation behaviour, but it cannot be used to select an architecture from protected final outcomes. The report consequently labels historical, development, and final evidence separately throughout.
